% =============================================================================
% PMVB Figure Legend
% =============================================================================
% Decoder for the box and wire conventions used across all TikZ-rendered
% figures in PMVB design documents (system block diagrams, functional block
% diagrams, schematics, power trees).
%
% Sequence diagrams (Mermaid) and timing diagrams (WaveDrom) render in their
% native styles and are documented separately.
%
% Compile:  pdflatex 00_legend.tex
% Convert:  pdftocairo -svg 00_legend.pdf
% =============================================================================
\documentclass[border=8pt]{standalone}
%
% =============================================================================
% pmvb-figures.sty -- Poor Man's Validation Bench figure style sheet
% =============================================================================
% Defines the house style for all TikZ figures in PMVB design documents.
%
% Usage in figure source:
%     \documentclass[border=4pt]{standalone}
%     \usepackage{pmvb-figures}
%
% Two color modes:
%     \pmvbDarkMode   -- dark background (default; matches SDD HTML theme)
%     \pmvbLightMode  -- white background (for printed portfolio PDFs)
% =============================================================================
\NeedsTeXFormat{LaTeX2e}
\ProvidesPackage{pmvb-figures}[2026/05/06 PMVB figure styles]
%
% --- Required packages ------------------------------------------------------
\RequirePackage{tikz}
% NOTE: dropped the siunitx option. Our figures use \pmvblabel{} for unit
% annotations rather than circuitikz's SI value formatting, and siunitx.sty
% isn't always available on bare TeX Live installs.
\RequirePackage[american]{circuitikz}
\RequirePackage{xcolor}
\RequirePackage{etoolbox}
%
% --- TikZ libraries ---------------------------------------------------------
\usetikzlibrary{arrows.meta}
\usetikzlibrary{positioning}
\usetikzlibrary{shapes.geometric}
\usetikzlibrary{shapes.misc}
\usetikzlibrary{fit}
\usetikzlibrary{backgrounds}
\usetikzlibrary{calc}
\usetikzlibrary{decorations.pathreplacing}
\usetikzlibrary{patterns}
%
% =============================================================================
% Color palette (FMCW dark theme, matches SDD HTML)
% =============================================================================
\definecolor{pmvbBg}{HTML}{0d1117}
\definecolor{pmvbFg}{HTML}{e6edf3}
\definecolor{pmvbMuted}{HTML}{94a3b8}
\definecolor{pmvbBlue}{HTML}{3b82f6}
\definecolor{pmvbBlueFill}{HTML}{1f2937}
\definecolor{pmvbGreen}{HTML}{10b981}
\definecolor{pmvbGreenFill}{HTML}{1e293b}
\definecolor{pmvbAmber}{HTML}{f59e0b}
\definecolor{pmvbRed}{HTML}{ef4444}
\definecolor{pmvbViolet}{HTML}{a78bfa}
%
\definecolor{pmvbBgLight}{HTML}{ffffff}
\definecolor{pmvbFgLight}{HTML}{0f172a}
\definecolor{pmvbMutedLight}{HTML}{475569}
\definecolor{pmvbBlueFillLight}{HTML}{dbeafe}
\definecolor{pmvbGreenFillLight}{HTML}{d1fae5}
%
% =============================================================================
% Mode selectors
% =============================================================================
\newcommand{\pmvbDarkMode}{%
  \pagecolor{pmvbBg}%
  \color{pmvbFg}%
}
%
\newcommand{\pmvbLightMode}{%
  \colorlet{pmvbBg}{pmvbBgLight}%
  \colorlet{pmvbFg}{pmvbFgLight}%
  \colorlet{pmvbMuted}{pmvbMutedLight}%
  \colorlet{pmvbBlueFill}{pmvbBlueFillLight}%
  \colorlet{pmvbGreenFill}{pmvbGreenFillLight}%
  \pagecolor{pmvbBgLight}%
  \color{pmvbFgLight}%
}
%
\AtBeginDocument{\pmvbDarkMode}
%
% =============================================================================
% Node styles
% -----------------------------------------------------------------------------
% IMPORTANT: each \tikzset{...} block below must be a single argument with NO
% blank lines inside, otherwise \par is inserted and pgfkeys errors out.
% =============================================================================
%
% --- Subsystem (system-level block: a board, a chip, a whole module) ---
\tikzset{
  pmvb subsystem/.style={
    rectangle,
    rounded corners=2pt,
    draw=pmvbBlue,
    line width=0.8pt,
    fill=pmvbBlueFill,
    text=pmvbFg,
    minimum width=28mm,
    minimum height=12mm,
    align=center,
    font=\sffamily\small,
    inner sep=4pt,
  },
}
%
% --- IC / hardware part (functional-level block) ---
\tikzset{
  pmvb ic/.style={
    rectangle,
    rounded corners=1pt,
    draw=pmvbGreen,
    line width=0.6pt,
    fill=pmvbGreenFill,
    text=pmvbFg,
    minimum width=24mm,
    minimum height=10mm,
    align=center,
    font=\sffamily\footnotesize,
    inner sep=3pt,
  },
}
%
% --- Internal block (inside an IC's functional block diagram) ---
\tikzset{
  pmvb internal/.style={
    rectangle,
    draw=pmvbFg,
    line width=0.5pt,
    fill=pmvbBg,
    text=pmvbFg,
    minimum width=18mm,
    minimum height=7mm,
    align=center,
    font=\sffamily\scriptsize,
    inner sep=2pt,
  },
}
%
% --- Pin pad (the box-with-X at IC die boundary in datasheet diagrams) ---
\tikzset{
  pmvb pin/.style={
    rectangle,
    draw=pmvbFg,
    line width=0.4pt,
    fill=pmvbBg,
    minimum size=2.5mm,
    inner sep=0pt,
    path picture={
      \draw[pmvbFg, line width=0.3pt]
        (path picture bounding box.south west) -- (path picture bounding box.north east)
        (path picture bounding box.north west) -- (path picture bounding box.south east);
    },
  },
}
%
% --- Connector (front panel, BNC, banana jack) ---
\tikzset{
  pmvb connector/.style={
    circle,
    draw=pmvbFg,
    line width=0.6pt,
    fill=pmvbBg,
    text=pmvbFg,
    minimum size=8mm,
    font=\sffamily\scriptsize,
    inner sep=1pt,
  },
}
%
% --- Annotation (no border, small muted italic text) ---
\tikzset{
  pmvb annotation/.style={
    text=pmvbMuted,
    font=\sffamily\scriptsize\itshape,
    align=center,
  },
}
%
% --- Subsystem grouping box (used with \node[fit=...]) ---
\tikzset{
  pmvb group/.style={
    rectangle,
    rounded corners=4pt,
    draw=pmvbMuted,
    line width=0.4pt,
    dashed,
    text=pmvbMuted,
    inner sep=8pt,
    font=\sffamily\scriptsize,
  },
}
%
% =============================================================================
% Edge styles
% -----------------------------------------------------------------------------
% Edge styles encode the *kind* of signal flowing on the line.
% =============================================================================
%
\tikzset{
  pmvb digital/.style={
    -Stealth,
    draw=pmvbBlue,
    line width=0.8pt,
    font=\sffamily\scriptsize,
  },
}
%
\tikzset{
  pmvb analog/.style={
    -Stealth,
    draw=pmvbAmber,
    line width=0.8pt,
    font=\sffamily\scriptsize,
  },
}
%
\tikzset{
  pmvb power/.style={
    -Stealth,
    draw=pmvbRed,
    line width=1.0pt,
    font=\sffamily\scriptsize,
  },
}
%
\tikzset{
  pmvb control/.style={
    -Stealth,
    draw=pmvbViolet,
    line width=0.7pt,
    dashed,
    font=\sffamily\scriptsize,
  },
}
%
\tikzset{
  pmvb bus/.style={
    -Stealth,
    draw=pmvbBlue,
    line width=1.4pt,
    double=pmvbBg,
    double distance=0.6pt,
    font=\sffamily\scriptsize,
  },
}
%
% Style for the "xN" bus-width annotation (overlaid on a bus line).
% Dark fill masks the underlying double-line so the white text reads cleanly.
\tikzset{
  pmvb bus width/.style={
    fill=pmvbBg,
    rounded corners=1pt,
    inner sep=2pt,
    font=\sffamily\bfseries\scriptsize,
    text=pmvbFg,
  },
}
%
\tikzset{
  pmvb optional/.style={
    -Stealth,
    draw=pmvbMuted,
    line width=0.6pt,
    dashed,
    font=\sffamily\scriptsize,
  },
}
%
% --- Bidirectional variants ---
\tikzset{
  pmvb digital bi/.style={
    pmvb digital,
    Stealth-Stealth,
  },
}
%
\tikzset{
  pmvb analog bi/.style={
    pmvb analog,
    Stealth-Stealth,
  },
}
%
% =============================================================================
% Diagram-level convenience
% =============================================================================
\tikzset{
  pmvb figure/.style={
    >=Stealth,
    every node/.append style={font=\sffamily},
    node distance=10mm and 14mm,
  },
}
%
% =============================================================================
% Helper macros
% =============================================================================
\newcommand{\pmvblabel}[1]{\textcolor{pmvbMuted}{\sffamily\scriptsize\itshape #1}}
%
\newcommand{\pmvbsubsystem}[3]{%
  node[pmvb subsystem] (#1) {{\bfseries #2}\\\scriptsize #3}%
}
%
\newcommand{\pmvbic}[3]{%
  node[pmvb ic] (#1) {{\bfseries #2}\\\scriptsize #3}%
}
%
% NOTE on bus-width annotations: TikZ's path parser doesn't reliably expand
% \newcommand macros at mid-path positions, so for the "xN" overlay use the
% inline node-in-path syntax directly:
%     \draw[pmvb bus] (A) -- node[pmvb bus width]{$\times 16$} (B);
% The pmvb bus width style is defined above.
%
% -----------------------------------------------------------------------------
% \opamptri{name}{coord}{dir}
%   Draws an op-amp triangle. Apex points right (dir=1) or left (dir=-1).
%   Defines coordinates: <name>_p, <name>_n, <name>_out for wiring.
%   Place any label as a separate \node call after the macro.
%
%   The {coord} argument MUST include its own parentheses, e.g.
%       \opamptri{bufa}{(0.6, -0.5)}{1}
%       \opamptri{tl072}{($(m1e.center) + (4mm, -14mm)$)}{1}
%   This lets callers use either simple coords or calc expressions.
% -----------------------------------------------------------------------------
\newcommand{\opamptri}[3]{%
  \begin{scope}[shift={#2}, xscale=#3]
    \draw[pmvbFg, line width=0.6pt, fill=pmvbBg]
      (0, 0) -- (
%
\begin{document}
\begin{tikzpicture}[pmvb figure]
  %
  % =========================================================================
  % Constants for layout
  % =========================================================================
  % Left column x-coord: where visual examples are anchored (west)
  % Right column x-coord: where descriptive text starts (anchor west)
  \def\colA{0}        % left column: visual element
  \def\colB{4.5}      % right column: description text
  \def\rowstep{0.9}   % vertical spacing between rows (cm)
  %
  % =========================================================================
  % Title bar
  % =========================================================================
  \node[font=\sffamily\bfseries, text=pmvbFg, anchor=west]
    at (\colA, 0) {Figure Legend};
  \node[font=\sffamily\scriptsize\itshape, text=pmvbMuted, anchor=west]
    at (\colA, -0.45)
    {Decoder for box and wire conventions used in PMVB block diagrams and schematics};
  %
  \draw[pmvbMuted, line width=0.4pt]
    (\colA, -0.85) -- (\colB + 9, -0.85);
  %
  % =========================================================================
  % Section: Boxes & Shapes
  % =========================================================================
  \node[font=\sffamily\bfseries\small, text=pmvbFg, anchor=west]
    at (\colA, -1.3) {Boxes \& Shapes};
  %
  % --- Row 1: Subsystem ---
  \node[pmvb subsystem, minimum width=24mm, minimum height=10mm, anchor=west]
    at (\colA, -2.0) {};
  \node[font=\sffamily\scriptsize, text=pmvbFg, anchor=west, align=left,
        text width=10cm]
    at (\colB, -2.0)
    {\textbf{Subsystem} -- a physical unit at platform level: a board (Pi 5), a
     piece of bench equipment (USB hub), or an entire module envelope.};
  %
  % --- Row 2: IC ---
  \node[pmvb ic, minimum width=24mm, minimum height=10mm, anchor=west]
    at (\colA, -3.0) {};
  \node[font=\sffamily\scriptsize, text=pmvbFg, anchor=west, align=left,
        text width=10cm]
    at (\colB, -3.0)
    {\textbf{IC / hardware part} -- a single chip or discrete part: an MCU,
     a DAC, a relay, an op-amp package, an FPGA.};
  %
  % --- Row 3: Internal block ---
  \node[pmvb internal, minimum width=24mm, minimum height=10mm, anchor=west]
    at (\colA, -4.0) {};
  \node[font=\sffamily\scriptsize, text=pmvbFg, anchor=west, align=left,
        text width=10cm]
    at (\colB, -4.0)
    {\textbf{Internal functional block} -- a stage inside an IC's functional
     block diagram (Interface Logic, DAC Register, String DAC, Output Logic).};
  %
  % --- Row 4: Connector ---
  \node[pmvb connector, anchor=west]
    at (\colA + 0.7, -5.0) {BNC};
  \node[font=\sffamily\scriptsize, text=pmvbFg, anchor=west, align=left,
        text width=10cm]
    at (\colB, -5.0)
    {\textbf{Connector} -- a physical input or output port: BNC, banana jack,
     USB-C, IEC inlet, terminal block.};
  %
  % --- Row 5: Op-amp triangle ---
  \opamptri{leg_oa}{(\colA + 1.5, -6.0)}{1}
  \node[font=\sffamily\scriptsize, text=pmvbFg, anchor=west, align=left,
        text width=10cm]
    at (\colB, -6.0)
    {\textbf{Op-amp / amplifier triangle} -- analog amplifier stage (in-IC
     or discrete). Coordinates \texttt{\_p}, \texttt{\_n}, \texttt{\_out}
     name the +/- inputs and the apex output.};
  %
  % =========================================================================
  % Section: Wires & Edges
  % =========================================================================
  \node[font=\sffamily\bfseries\small, text=pmvbFg, anchor=west]
    at (\colA, -7.3) {Wires \& Edges};
  %
  % --- Row 6: Digital signal ---
  \draw[pmvb digital] (\colA, -8.0) -- ++(2.4, 0);
  \node[font=\sffamily\scriptsize, text=pmvbFg, anchor=west, align=left,
        text width=10cm]
    at (\colB, -8.0)
    {\textbf{Digital signal} -- single-bit logic line: SPI MOSI/MISO/SCK,
     GPIO, USB-TMC command/data, $\overline{\text{CS}}$.};
  %
  % --- Row 7: Multi-bit bus, with width annotation ---
  % Macro \pmvbbus{N} doesn't expand cleanly into a TikZ path's mid-position,
  % so we inline the node-in-path here directly. (The macro is still useful
  % when its expansion happens at a syntactic boundary other than mid-path.)
  \draw[pmvb bus] (0, -8.9) -- node[pmvb bus width]{$\times 16$} (2.4, -8.9);
  \node[font=\sffamily\scriptsize, text=pmvbFg, anchor=west, align=left,
        text width=10cm]
    at (\colB, -8.9)
    {\textbf{Multi-bit bus} -- aggregated parallel link with optional
     ``$\times N$'' width annotation overlaid on the line (USB-TMC fabric,
     parallel ADC/DAC bus, FPGA capture stream).};
  %
  % --- Row 8: Analog signal ---
  \draw[pmvb analog] (\colA, -9.8) -- ++(2.4, 0);
  \node[font=\sffamily\scriptsize, text=pmvbFg, anchor=west, align=left,
        text width=10cm]
    at (\colB, -9.8)
    {\textbf{Analog signal} -- continuous voltage or current path: DAC output,
     op-amp output, sensor input, BNC analog stimulus.};
  %
  % --- Row 9: Power rail (no arrow head) ---
  \draw[pmvbFg, line width=1.0pt, draw=pmvbRed] (\colA, -10.7) -- ++(2.4, 0);
  \node[font=\sffamily\scriptsize, text=pmvbFg, anchor=west, align=left,
        text width=10cm]
    at (\colB, -10.7)
    {\textbf{Power rail} -- supply voltage line. No arrow head (passive
     distribution, not a directional signal). Use red for $V_{DD}$/+5V/+12V,
     and a ground glyph for $V_{SS}$/GND.};
  %
  % --- Row 10: Control signal (dashed violet) ---
  \draw[pmvb control] (\colA, -11.6) -- ++(2.4, 0);
  \node[font=\sffamily\scriptsize, text=pmvbFg, anchor=west, align=left,
        text width=10cm]
    at (\colB, -11.6)
    {\textbf{Control / strobe} -- synchronous control or load-enable line:
     $\overline{\text{LDAC}}$, $\overline{\text{SHDN}}$, trigger arm, latch
     enable.};
  %
  % --- Row 11: Optional / deferred (dashed gray) ---
  \draw[pmvb optional] (\colA, -12.5) -- ++(2.4, 0);
  \node[font=\sffamily\scriptsize, text=pmvbFg, anchor=west, align=left,
        text width=10cm]
    at (\colB, -12.5)
    {\textbf{Optional / deferred} -- conditional or future connection:
     shared trigger bus, v2.0 streaming sidecar link, planned upgrade path.};
  %
  % =========================================================================
  % Section: Text styles
  % =========================================================================
  \node[font=\sffamily\bfseries\small, text=pmvbFg, anchor=west]
    at (\colA, -13.8) {Text Styles};
  %
  % --- Row 12: Primary label ---
  \node[font=\sffamily\bfseries\scriptsize, text=pmvbFg, anchor=west]
    at (\colA, -14.5) {Pico 2 W};
  \node[font=\sffamily\scriptsize, text=pmvbFg, anchor=west, align=left,
        text width=10cm]
    at (\colB, -14.5)
    {\textbf{Primary label} (off-white, bold or regular sans) -- subsystem
     names, IC names, pin labels.};
  %
  % --- Row 13: Wire annotation ---
  \node[font=\sffamily\scriptsize\itshape, text=pmvbMuted, anchor=west]
    at (\colA, -15.4) {SCPI commands + data};
  \node[font=\sffamily\scriptsize, text=pmvbFg, anchor=west, align=left,
        text width=10cm]
    at (\colB, -15.4)
    {\textbf{Wire annotation} (gray italic, via \texttt{\textbackslash pmvblabel})
     -- explanatory text on edges: bus names, voltage ranges, frequencies,
     bit widths.};
  %
  % --- Row 14: Power label ---
  \node[font=\sffamily\scriptsize, text=pmvbRed, anchor=west]
    at (\colA, -16.3) {+3.3 V};
  \node[font=\sffamily\scriptsize, text=pmvbFg, anchor=west, align=left,
        text width=10cm]
    at (\colB, -16.3)
    {\textbf{Power label} (red text) -- voltage rail identifiers placed
     adjacent to a power node or rail.};
  %
\end{tikzpicture}
\end{document}
